\documentclass[11pt,a4paper]{article}
\usepackage[T1]{fontenc}
\usepackage[utf8]{inputenc}
\usepackage{lmodern}
\usepackage[margin=28mm]{geometry}
\usepackage{amsmath,amssymb,amsthm,mathtools}
\usepackage{booktabs,longtable,array}
\usepackage{algorithm,algpseudocode}
\usepackage[section]{placeins}
\usepackage{needspace}
\floatstyle{boxed}
\restylefloat{algorithm}
\usepackage{xcolor}
\usepackage[framemethod=tikz]{mdframed}
\usepackage{microtype}
\usepackage{xurl}
\usepackage[hidelinks]{hyperref}
\hypersetup{
  pdftitle={Free FT: Ternary Lattice Signatures with a Code-Bound Security Reduction},
  pdfauthor={Niirmata},
  pdfsubject={Revised research draft 0.3; conditional classical reduction; ePrint preparation},
  pdfkeywords={FT1536, Falcon, EUF-CMA, MT-ISIS, Lean, source binding}
}
\theoremstyle{plain}
\newtheorem{theorem}{Theorem}[section]
\newtheorem{proposition}[theorem]{Proposition}
\newtheorem{lemma}[theorem]{Lemma}
\newtheorem{corollary}[theorem]{Corollary}
\theoremstyle{definition}
\newtheorem{definition}[theorem]{Definition}
\newtheorem{assumption}[theorem]{Assumption}
\newtheorem{target}[theorem]{Target statement (in flight)}
\theoremstyle{remark}
\newtheorem{remark}[theorem]{Remark}
\newcommand{\CLOSED}{\textsf{CLOSED}}
\newcommand{\OPEN}{\textsf{OPEN}}
\newcommand{\INFLIGHT}{\textsf{IN-FLIGHT}}
\newcommand{\UNCERTIFIED}{\textsf{UNCERTIFIED}}
\newcommand{\NOTREVIEWED}{\textsf{NOT\_REVIEWED}}
\newcommand{\ALTP}{\textsf{ALTERNATIVE\_PROPOSAL}}
\newcommand{\status}[1]{\par\noindent{\small\textbf{Scope.} #1.}\par\smallskip}
\newcommand{\todo}[2]{\par\smallskip\noindent\textbf{\OPEN{} [#1].} #2\par\smallskip}
\newcommand{\src}[1]{\textsuperscript{\textsf{[#1]}}}
\newcommand{\code}[1]{\texttt{\detokenize{#1}}}
\newcommand{\Adv}{\operatorname{Adv}}
\newcommand{\Law}{\operatorname{Law}}
\newcommand{\Second}{\operatorname{second}}
\newcommand{\AC}{\operatorname{AC}}
\newcommand{\TV}{\operatorname{TV}}
\newcommand{\none}{\bot}
\newcommand{\Rq}{R_q}
\newcommand{\epscoll}{\varepsilon_{\mathrm{coll}}}
\newcommand{\dprg}{\delta_{\mathrm{PRG}}}
\newcommand{\Z}{\mathbb{Z}}
\newcommand{\R}{\mathbb{R}}
\newcommand{\U}{\mathsf{U}}
\newcommand{\ind}{\mathbf{1}}
\AddToHook{cmd/section/before}{\Needspace{8\baselineskip}}
\input{build/snapshot}
\title{Free FT: Ternary Lattice Signatures\\
  with a Code-Bound Security Reduction}
\author{Niirmata\\\small\url{https://github.com/niirmataa/free_falcon_sign}}
\date{Research draft 0.3, revised 10 October 2026\\\small Prepared for ePrint}
\begin{document}
\maketitle
\begin{abstract}
FT1536 is a ternary-secret signature profile over a degree-$1536$ NTRU
ring with modulus $18433$. We present its conditional classical
EUF-CMA-to-MT-ISIS reduction and the source-bound components being
developed to instantiate that reduction. The formal game supplies a
concrete reducer, an explicit collision loss, and a directional
second-moment comparison. Under named pointwise and law premises, the
sampler analysis computes $e_2(k)=k^{32}-1<2^{-32}$ for an unconditioned
attempt shape; an acceptance-conditioned shape instead gives
$e_2(k_{\mathrm{cond}})<2^{-17}$. Identifying the realized shape remains an obligation.
A separate computational theorem treats public simulation using a
restricted class of randomized tests and an explicit cost certificate.
Source-bound execution results establish exact integer NTRU and public
inverse equations at stated fragment boundaries. Kernel counterexamples
rule out bulk-only, total-mass and additive-error shortcuts, and exclude
a small statistical error for long tapes from the modeled finite-state
generator. The binding process also exposed a latent C buffer-cursor
defect and defects in the formal parsers. An evidence ledger distinguishes
these results from finite validation and unresolved interfaces. Lattice
estimates remain uncertified S06 diagnostics; symmetric search ceilings
are reported separately, with system accounting given by the minimum.
The honest-signer bridge, emitted-key law and global byte-level
composition remain open. The paper claims a conditional classical
reduction, not end-to-end implementation security or a QROM reduction.
\src{C01,C03,C04,C06,C08,C09,C10,C12,C14,C15,C20,C21}
\end{abstract}
\medskip
\noindent{\small\textbf{Evidence snapshot.} SHA-256 of \code{SOURCES.sha256}:\\
\texttt{\EvidenceManifestSHA}\\
\textbf{Editorial base.} Git \texttt{\EditorialBaseCommit}.}
\par\smallskip
\noindent\textbf{Reading convention.} \CLOSED{} labels refer to the stated
formal component, including its premises, rather than to implementation
security or independent acceptance. \INFLIGHT{} and \OPEN{} labels identify
undischarged work. Superscript identifiers refer to the exact paths,
declarations and byte pins in the companion \code{SOURCES.md}; every
number and quoted statement is traceable through that map. Statements
printed as targets express the intended assembly and are not theorems.
Full snapshot bindings and reproduction instructions are in
Appendix~\ref{app:trust}.
\clearpage
\section{Introduction and purpose}\label{sec:intro}
\status{CLOSED narrative material; owner voice review pending}
The objective of the FT1536 development is a security argument whose
objects can be traced to a specified implementation. Such an argument
requires more than a reduction between ideal games: the key law,
sampler law, randomness interface and byte-level acceptance predicate
must refer to the same experiment. The present artifact contains a
conditional classical reduction and several source-bound components;
their final implementation-level composition is unfinished.\src{C01}

This paper develops three parts of that argument. First, it presents the
conditional reduction and the computed constants of a pointwise sampler
comparison. Second, it records negative results that rule out shortcuts
in that comparison. Third, it describes the source-binding and review
process, including defects in both C code and formal models that the
process exposed. The artifact is part of the FT1536 evidence
archive~\cite{ft1536artifact}.\src{C04,C06,C09,C11,C12}

% Mission paragraph signed by the owner 2026-10-07 ("idealna wersja do
% paperu i do serca") - EN text below, PL original recorded in
% notes/PAPER_WORK_STATE.md.
FT1536 exists because trust in signature software should rest on
arguments that anyone can inspect, not on the authority of whoever
shipped the binary. The project develops ternary-secret signature
profiles in the Falcon lineage --- freely chosen, freely auditable,
freely usable --- to make independently inspectable cryptographic tools
available to people who wish to choose and evaluate such tools
themselves. That purpose requires publishing limitations and failed
arguments alongside successful proofs, binding claims to the exact code
that runs, and preserving the attribution of the Falcon Project and
Thomas Pornin. Accessibility is a design aim; it is not evidence for a
security claim.

Section~\ref{sec:main} distinguishes the target assembly from the
theorems already present. Section~\ref{sec:sampler} states the negative
results as theorems. Sections~\ref{sec:security} and~\ref{sec:limits}
give the security-accounting convention and the unresolved obligations
in the body of the paper.
The end-to-end assembly and its in-flight obligations are described in
Section~\ref{sec:main}.

\section{Preliminaries and threat model}\label{sec:model}
\status{CLOSED definitions; source instantiation IN-FLIGHT}
\subsection{The relation and its geometry}
The formal representation of $\Rq$ is
$\operatorname{Fin}(768)\to(\Z/18433\Z)^2$.
Multiplication is specified by polynomial remainder modulo
$X^{1536}-X^{768}+1$, not by an NTT oracle. Write
$V=\operatorname{Fin}(768)\to\Z^2$, and define
\begin{align*}
 Q_0(v)&=\sum_{i=0}^{767}\bigl(v_{i,1}^2+v_{i,1}v_{i,2}+v_{i,2}^2\bigr),\\
 Q(u,v)&=Q_0(u)+Q_0(v),\qquad B=2093922385,\\
 A_h(u,v)&=\operatorname{reduce}(u)+h\operatorname{reduce}(v).
\end{align*}
\code{ShortPreimage} is exactly $A_h(u,v)=c\ \wedge\ Q(u,v)<B$.
The coefficient-valued verifier accepts $s$ precisely when it is signed-16
and
$Q(\operatorname{center}(c-h\operatorname{reduce}(s)),s)<B$.
The centering interval is $[-9216,9216]$.\src{C03}

\begin{proposition}[Accepted coefficients extract a short preimage]
For every $h,c\in\Rq$ and $s\in V$, if the verifier accepts $s$ for
$(h,c)$, then the extracted pair is a short preimage for $(h,c)$.
\end{proposition}
\begin{proof}
\code{reduce_center} gives the equation for the extracted pair;
the norm inequality is the second conjunct of \code{Verify}.
This is \code{Relation.accepted_extracts}.\src{C03}
\end{proof}

In the MT-ISIS experiment, $h$ is drawn from a specified public-key law
$\mu_H$, and the adversary receives $T$ independent uniform targets
in $\Rq$. It returns an index and a pair; the experiment checks the
index and the same \code{ShortPreimage} relation. Hardness must therefore
refer to this key law, target law, quadratic form and resource class.
An unrelated homogeneous-SIS estimate is not a replacement.\src{C05}

\subsection{Classical interactive games}
A budget record $\beta$ contains $q_s,q_h$, an adversary coin length and
a byte parameter. The type \code{Program} enforces the Sign and hash
query counts structurally: a signing query consumes a token even when
the reply is a failure. The adversary receives the public key and its
coins. Hash names are parsed into a $40$-byte nonce and a message when
long enough, with a separate short-input case. Successful forgery
requires a message absent from the submitted-message table and a nonce
of the specified length. The proof does not turn the byte-budget field
alone into a calibrated runtime bound.\src{C05}

All oracle queries and interactions here are classical. The formal
\code{UniformChallengeAt S} interface requires the challenge marginal
of the sampler law to equal the uniform law at each public key, history,
message and nonce. It names an exact ROM interface, not a proof that
the concrete SHAKE implementation is a random oracle. No QROM claim is
made.\src{C01,C18}

\subsection{Conditioning and operational assumptions}
Decision D1 specifies the emitted key law as
$\Law(\text{attempt}\mid\mathrm{Accept})$, with public marginal $\mu_H$.
In the IID per-attempt model with acceptance probability $p>0$, a cap
$L$ has availability $1-(1-p)^L$ and the first successful attempt,
conditioned on a success, has that conditional law. These are modeling
identities recorded by D1/S1; binding the real restarted KeyGen to them
is a B1 obligation. The empirical reciprocal in
Section~\ref{sec:parameters} must not be substituted as a proved
conditioning factor. Key conditioning specifies the law in the MT
assumption; it is not an extra automatic multiplier of the Sign
sampler's $e$.\src{C01,C02}

Decision D2 selects a computational change of randomness, using the
restricted-test premise of Section~\ref{sec:computational}. Decision
D3 requires a proved sampler constant at the stated interfaces.
We retain the named operational boundaries:
\begin{description}
 \item[A1.] KeyGen timing and restart counts are absent from the
 adversary's view. This is a deployment assumption, not a timing proof.
 \item[A2.] Seed entropy and independence are idealized at the seed
 boundary. ROM uniformity is separately named at the challenge boundary;
 neither premise proves the other.
 \item[A3/A4.] The implementation verdict and serialization must match
 the formal relation and framing. Their final byte-level composition
 remains open.
 \item[A5.] Source contracts use the stated LP64 word and wrapping
 semantics. Faults and side channels are outside this game.
\end{description}
The complete ledger appears in Section~\ref{sec:limits}.\src{C01}

\subsection{Directional second moment}
For finite laws $j,p$, the formal definitions are
\[
 \AC(j,p)\iff\forall x,\ p(x)=0\Rightarrow j(x)=0,
 \qquad \Second(j,p)=\sum_x\frac{j(x)^2}{p(x)}.
\]
Under absolute continuity, this is $1+\chi^2(j\Vert p)$.
The separate absolute-continuity premise is essential: Lean's real
division is total, while the extended-valued \code{chi2} definition is
infinite when absolute continuity fails. The direction used below is
simulation relative to the honest law, not the reverse.\src{C03,C04}

\section{The FT family and parameters}\label{sec:parameters}
\status{CLOSED profile facts and measurement record; family certification IN-FLIGHT}
\todo{DRAFT-02}{Transcribe the pinned profile, acceptance measurement and codec scope.}

\section{Main statement and assembly interfaces}\label{sec:main}
\status{CLOSED conditional components; real-Sign assembly IN-FLIGHT}
\todo{B1-HONEST-HOP}{Bind a tape-driven honest signer and its certified winning test.}
\todo{B1-LAWS}{Supply the fair-tape honest-game binding, realized sampler certificate and key-law identification.}

\section{The classical reduction}\label{sec:reduction}
\status{CLOSED conditional reduction}
\todo{DRAFT-05}{Describe the concrete reducer and directional second-moment comparison.}

\section{Sampler analysis and negative results}\label{sec:sampler}
\status{CLOSED arithmetic and conditional transports; realized weights IN-FLIGHT}
\todo{DRAFT-06}{State both attempt shapes and the kernel counterexamples with exact hypotheses.}

\section{Binding to the implementation}\label{sec:binding}
\status{CLOSED method and recorded findings; B1 IN-FLIGHT}
\subsection{From source text to executions}
The source3 method fixes C bytes, parses the supported C99 fragment into
a statement tree, and proves the equality between the parsed region
and that tree. Execution theorems then reason about this bound tree
under explicit state, memory and word-range contracts. This separates
the source identity from the value-level argument. It is a fragment
semantics, not a verified general C compiler.\src{C11}

For example, \code{KeygenMkgm3Program.source_bound} states
\[
 \code{C99ModularParser.region}\ 2945\ 91=\mathrm{some}\ \code{code}.
\]
The B1.02 execution work derives loop counters and pointer positions
for the first, triple and intermediate loop families of
\code{modp_NTT3_ext}. Butterfly-call observations are conclusions of
executions of the parsed bodies; assuming the intended calls would
leave the source-binding obligation unanswered. These are scoped
execution results, still subject to their memory and non-overflow
premises.\src{C11}

The generator certificate is similarly concrete. With
$p=2147355649$ and $g=1907584673$, the declarations in
\code{KeygenGeneratorOrder} prove
\begin{align*}
 g^{9216}\bmod p&=1,&g^{4608}\bmod p&\ne1,&g^{3072}\bmod p&\ne1,\\
 (g^2)^{4608}\bmod p&=1,&(g^2)^{2304}\bmod p&\ne1,&
 (g^2)^{1536}\bmod p&\ne1.
\end{align*}
The same module checks $9216=2^{10}3^2$ and $4608=2^9 3^2$.
These are the exact power tests used as the order certificates for
$9216$ and $4608$; they do not assume an NTT correctness oracle.
We quote the literal arithmetic declarations rather than silently
replacing them by a stronger source-level transform theorem.\src{C11}

\subsection{The verification found bugs}
\paragraph{F-001: an inherited C buffer-cursor defect.}
In the pinned \code{falcon_prng_get_bytes}, the available length is
computed from \code{p->ptr}, but the copy is
\code{memcpy(buf, p->buf.d, clen)} rather than a copy from
\code{p->buf.d + p->ptr}. A call beginning at a nonzero cursor can
therefore copy the buffer prefix instead of the intended stream slice.
The recorded census found the same statement in seven vendored
reference/optimized/local copies and no call sites in the examined
profiles. This is an inherited defect in the upstream lineage as
vendored here, classified LATENT; it is not evidence about every later
upstream release. The examined signing path uses the indexed
\code{get_u8}/\code{get_u64} accessors. A future use of the bulk-byte
API would activate the faulty path.\src{C12}

\paragraph{A dead REV10 parser branch.}
The initial store parser never recognized real expressions of the form
\code{gm[b + REV10[u << k]]}. Its greedy \code{pureExpr} consumed
\code{+ REV10} as part of a scalar expression, leaving a bracket behind.
The correction splits at the REV10 marker before parsing the base.
An apparently plausible row model would otherwise have described an
unbound program tree. This was a formal-model defect, not a new defect
in the C table.\src{C12}

\paragraph{Duplicated XOR and a false function call.}
The candidate \code{divSelect} included an outer XOR with \code{z},
although the compound update \code{z ^= ...} already supplied that
operation. Kernel checking rejected the claimed parse equality.
The correction uses precisely the right-hand side
\code{(z ^ z2) & -(uint32_t)((e >> i) & 1)}.
A separate parser pattern treated an opening parenthesis as a function
name and misread \code{((size_t)1 << k) - 1}; a token-name guard removed
the false call. The failed bindings and probes are retained in the
batch record.\src{C12}

\subsection{Proof engineering and current boundary}
The REV10 certificate now states
\code{rawTable = some ((List.range 1024).map bitrev10)}.
Its source-side certificate uses eleven slices, each parsing at most
eight pinned lines, and composes intermediate equalities without
re-parsing the entire table. Eight lines was the successful granularity
in the recorded experiment, not a universal limit of Lean's kernel.
The final named job is \code{keygen_rev10_cert_004}; the checkpoint
records a clean accepted run.\src{C11}

These results illustrate why a kernel proof of a model and a proof
that the model describes the source are distinct obligations. The
remaining B1 work includes source value/row laws and the eventual
emitted-material-to-fiber and law interfaces. Their completion cannot
be inferred from the closed parser or table certificates.\src{C11,C01}
\todo{B1-FINAL}{Replace this progress stub with the final source-to-law exports after delivery.}

\section{Security accounting: lattice estimates and symmetric ceilings}\label{sec:security}
\status{CLOSED diagnostic record; S06 certification OPEN}
\todo{DRAFT-08}{Present model-dependent lattice estimates and symmetric ceilings separately.}

\section{The computational interface}\label{sec:computational}
\status{CLOSED conditional public-simulation component; independent acceptance OPEN}
\todo{DRAFT-09}{Give the admitted-class hop, full-run tape and actual cost certificate.}

\section{Related work}\label{sec:related}
\status{OPEN literature review}
\todo{LITERATURE}{Complete the comparison with GPV, NTRU, BLISS, Falcon reductions, verified compilers and implementation verification.}

\section{Limitations and open work: the honesty ledger}\label{sec:limits}
\status{CLOSED scope record; listed proof obligations OPEN or IN-FLIGHT}
\todo{DRAFT-11}{Transcribe the A1--A5 and B1--B5 ledger, including the byte-support and QROM boundaries.}

\bibliographystyle{plain}
\bibliography{refs}
\clearpage
\appendix
\section{Trust chain and reproducibility}\label{app:trust}
\status{CLOSED artifact methodology; release packaging OPEN}
\todo{DRAFT-A}{Document pins, receipts, axiom audits and the actual independence of the reviews.}

\section{Project provenance}\label{app:origin}
\status{\OPEN{} historical documentation}
The project's stated purpose appears in Section~\ref{sec:intro}.
Chronology requires dated source material of its own; earlier drafts
and the current proof snapshot do not by themselves establish the
project's first creation date. This appendix reserves that record for
the owner-supplied material and its eventual artifact bindings.
\todo{ORIGIN}{The owner must supply the March-disk material and an authenticated origin timeline. The dates suggested in the outline are not yet established here.}

\end{document}
